\documentclass[manuscript,review,anonymous]{acmart}
\acmJournal{TIIS}
\setcounter{tocdepth}{3}
\usepackage{array}
\usepackage{booktabs} % removed duplicate
\usepackage{pdfpages}
\usepackage{longtable}
\usepackage{pdflscape}
\usepackage{tabularx}
\usepackage{xltabular}
\usepackage{ragged2e}
\usepackage{subcaption}
\usepackage{multirow}
\usepackage{graphicx}
\usepackage{dblfloatfix}
\usepackage{placeins}

% Table/Figure A1, A2 formatting
\renewcommand{\thefigure}{A\arabic{figure}}
\renewcommand{\thetable}{A\arabic{table}}
\setcounter{figure}{0}
\setcounter{table}{0}
\newcommand{\ts}{\textsuperscript}
\newcommand{\tsub}{\textsubscript}

\newcolumntype{A}[1]{>{\raggedright\arraybackslash}p{#1}}
\newcolumntype{L}{>{\raggedright\arraybackslash}X}
\usepackage[super]{nth}
\usepackage[nameinlink,noabbrev]{cleveref}

\AtBeginDocument{%
  \providecommand\BibTeX{{%
    Bib\TeX}}}
\setcopyright{acmlicensed}
\copyrightyear{2026}
\acmYear{2026}
\acmDOI{XXXXXXX.XXXXXXX}

\usepackage[acronym]{glossaries}
\input{acronyms} 

\begin{document}

% --- TITLE FIX HERE ---
\title{Online Appendix: [From Experts to Language Models: Enhancing Online Expert Elicitation with LLMs]}
% Optional subtitle to link back to the main submission:
% \subtitle{Supplementary Material for "Your Main Paper Title"}

\maketitle

\keywords{Expert elicitation, Large language models, Judgment errors}

\section{Document Preamble}
This document present evaluation of the use of seed, and shuffling of the list of \gls{dc} solutions and the barriers. The target is to check whether those changes may bring different results from the ones already presented in the manuscript. 

%\clearpage

\section{Experimental Design}
We prompted the models on feasibility and barrier questions. The two questions presentation to the \glspl{llm} involves the \gls{dc} solution and the barrier list presentationns in the prompts. Whereas, for the importance question each \gls{dc} solution was presented in isolation. 
For feasibility question we performed deterministic shuffling of the \gls{dc} solutions with seed (two iterations), and without seed (three runs).
As for the barrier-select as well we had the deterministic and the non-deterministic shuffling using Fisher-Yates algorithm mixing with and without a seed. In both cases  we used the previous response we had as a baselines of comparisons. 
Further in the barrier-select we fitted a regression model to assess the primacy and recency effect of the barrier list presentation using the randomized order responses. 
The list of eleven barriers was divived as follows: first three (for primacy), last three for recency, and the five in the middle for the middle effect.

\section{Results}

\subsection{Feasibility Question}
Gemma3:12B, and mistral did not alter their ratings in eaither of the experimental configurations. 
While for Llama-Pro and Phi-4 they exhibitted similar behavioral patterns. 
While their ratings resembles in on baseline to the three runs (without seed); whereas with a seed ragardless of the \gls{dc} solutions arrangements still generated similar responses. 
Llama-Pro had all ratings 2 (somewhat feasible) and Phi-4 had (67\%) on feasible, with (33\%) on very feasible. 

These results demonstrate the reproducibility of the responses for all of the four model variants when seed is included during prompting. 
On the other hand, shuffling of the barrier lists does not guarantee models variation in the responses. 
Therefore due to these findings we re-run the experiments from zero-shot, context, to anchor to include a seed. This was done for both importance and feasibility questions. 


\subsection{Barrier-select Question}



\end{document}